%!TEX TS-program = xelatex
%!TEX encoding = UTF-8 Unicode

\documentclass[11pt, a4paper]{article}
\usepackage{fontspec} 

% DOCUMENT LAYOUT
\usepackage{geometry} 
\geometry{a4paper, textwidth=5.5in, textheight=8.5in, marginparsep=7pt, marginparwidth=.6in}
\setlength\parindent{0in}
\emergencystretch=1.5em

% FONTS
\usepackage[usenames,dvipsnames]{xcolor}
\usepackage{xunicode}
\usepackage{xltxtra}
\defaultfontfeatures{Mapping=tex-text}
%\setromanfont [Ligatures={Common}, Numbers={OldStyle}, Variant=01]{Linux Libertine O}
%\setmonofont[Scale=0.8]{Monaco}
%%% modified by Karol Kozioł for ShareLaTeX use
\setmainfont[
  Path=../,
  Ligatures={Common}, Numbers={OldStyle}, Variant=01,
  BoldFont=LinLibertine_RB.otf,
  ItalicFont=LinLibertine_RI.otf,
  BoldItalicFont=LinLibertine_RBI.otf
]{LinLibertine_R.otf}
\setmonofont[Scale=0.9, Path=../]{DejaVuSansMono.ttf}

% ---- CUSTOM COMMANDS
\chardef\&="E050
\newcommand{\html}[1]{\href{#1}{\scriptsize\textsc{[html]}}}
\newcommand{\pdf}[1]{\href{#1}{\scriptsize\textsc{[pdf]}}}
\newcommand{\doi}[1]{\href{#1}{\scriptsize\textsc{[doi]}}}
% ---- MARGIN YEARS
\usepackage{marginnote}
\newcommand{\amper{}}{\chardef\amper="E0BD }
\newcommand{\years}[1]{\marginnote{\scriptsize #1}}
\renewcommand*{\raggedleftmarginnote}{}
\setlength{\marginparsep}{11pt}
\reversemarginpar

% HEADINGS
\usepackage{sectsty} 
\usepackage{multicol} 
\usepackage[normalem]{ulem} 
\sectionfont{\mdseries\upshape\Large}
\subsectionfont{\bfseries\upshape\normalsize} 
\subsubsectionfont{\mdseries\upshape\large} 

% PDF SETUP
% ---- FILL IN HERE THE DOC TITLE AND AUTHOR
\usepackage[%dvipdfm, 
bookmarks, colorlinks, breaklinks, 
% ---- FILL IN HERE THE TITLE AND AUTHOR
	pdftitle={Junjun Yin - CV (Russell Sage Foundation)},
	pdfauthor={Junjun Yin},
	pdfproducer={}
]{hyperref}  
\hypersetup{linkcolor=blue,citecolor=blue,filecolor=black,urlcolor=MidnightBlue} 

% DOCUMENT
\begin{document}
{\LARGE Junjun Yin, Ph.D}
\begin{center}
\rule{\textwidth}{0.4pt}
\end{center}
\section*{Contact Information}

\begin{multicols}{2}
 \begin{flushleft}
\emph{Assistant Professor}\\
Department of Data Science\\
Columbian College of Arts \& Sciences\\
The George Washington University\\
Washington, DC 20052
\end{flushleft}
\columnbreak
\begin{flushright}
Address: Room 309, Samson Hall\\
Phone: +1 217-819-6366\\[.1cm]
Email: \href{mailto:j.yin@gwu.edu}{j.yin@gwu.edu}\\[.1cm]
Web: \href{https://www.junjunyin.com}{www.junjunyin.com}
\end{flushright}
\end{multicols}
%%\hrule

\section*{Academic Appointments}
\noindent
\years{2025-}Assistant Professor, Department of Data Science, The George Washington University, USA\\
\years{2019-2025}Assistant Research Professor, Social Science Research Institute, Pennsylvania State University, USA\\
\years{2019-2025}ICDS Associate, Institute for Computational and Data Sciences, Pennsylvania State University, USA\\
\years{2016-2018}Social Science Big Data Research Scientist, Pennsylvania State University, USA\\
\years{2016} Lecturer, Department of Geography and Geographic Information Science; University of Illinois at Urbana-Champaign, USA\\
\years{2014-2016} Postdoctoral Research Fellow, CyberGIS Center for Advanced Digital and Spatial Studies; National Center for Supercomputing Applications; University of Illinois at Urbana-Champaign, USA\\
\years{2013-2014} Assistant Lecturer, Technological University Dublin, Ireland
%%\hrule

\section*{Education}
\noindent
\years{2009 - 2013}\textit{Ph.D} in Spatial Information Science, Technological University Dublin, Ireland\\
Science Foundation Ireland Scholarship\\
\textit{Advisors: Dr. James D. Carswell (Technological University Dublin), Dr. Michela Bertolotto (University College Dublin)}\\
\years{2006 - 2009}\textit{MSc} in Geoinformatics, University of Gävle, Sweden\\
\textit{Advisor: Dr. Bin Jiang (University of Gävle)}\\
\years{2002 - 2006}\textit{BSc} in Electronics Engineering, University of Electronic Science and Technology of China

\section*{Research Interests}
GIScience; Geospatial Data Science; Geographic Big Data;
Human-Urban Interactions; Human Mobility; Urban Informatics; (Geo)Visual-Analytics; 
Spatial Data Mining and Deep Learning (GeoAI); Spatial/Graph Neural Networks; High-performance (cyberGIS), Web, and Mobile GIS\\

%\hrule

\section*{Selected Publications}
\subsection*{Submitted}
\years{2026}\textbf{Yin, J.}, Yin, L., Du, J., Yang, D., Hu, B. and Jiang, B., Empirical Validation of Architectural Livingness: Testing Alexander’s Theory through AI-Based Measurement, Perceptual Response, and Brain Activity. (\textit{submitted to Journal of Environmental Psychology, in review})\\
\years{2026}Chi, G., \textbf{Yin, J.}, Van Hook,  J., Plutzer, E. and Xu, H., The Mis-representativeness of Twitter Data. (\textit{submitted to Demography, in review})\\
\years{2026}Son, J.H. and \textbf{Yin, J.}, Beyond Residence: Mobility-Driven PM2.5 Exposure Disparity in the Washington, D.C. Metropolitan Area. (\textit{submitted to Journal of Transport \& Health, in review})\\
\years{2026}Lakshmanan, G.S. and \textbf{Yin, J.}, AirGuard: A Machine Learning Framework for High-Resolution PM2.5 Prediction in California Using Spatiotemporal Data Integration. (\textit{submitted to Environmental Data Science, in review})

\subsection*{Peer Reviewed Journals}
\years{2026}Smith, M.L., Chi, G., \textbf{Yin, J.}, Bodovski, Y. and Sundar, S.S. (2026), County Partisanship Affects Social Media Posting Behavior During a Pandemic. \textit{Social Media + Society}, 12(1), pp. 1-12. (DOI: 10.1177/20563051261419387)\\
\years{2025}\textbf{Yin, J.} Unlocking the Secrets of Scenic Beauty:  An artificial intelligence approach to understanding public perceptions of environmental aesthetics. \textit{Journal of the Royal Society Interface}, 22(225): 20250045. (DOI: 10.1098/rsif.2025.0045)\\
\years{2025}Sundar, S.S., Snyder, E., Liao, M., \textbf{Yin, J.}, Wang, J., and Chi, G., Sharing without clicking: The concerning phenomenon of mindless news forwarding on social media. \textit{Nature Human Behaviour}, 9(1), pp. 156-168. (DOI:10.1038/s41562-024-02067-4)\\
\years{2023}\textbf{Yin, J.} and Chi, G. (2023), A Tale of Three Cities: Uncovering human-urban interactions with geographic-context aware social media data. \textit{Urban Informatics, 1(1), 20}\\(DOI:10.1007/s44212-022-00020-2)\\
\years{2023}Yu, M., Zhang, S., Zhang, K., \textbf{Yin, J.}, Varela, M. and Miao, J. (2023), Developing high-resolution PM2.5 exposure models by integrating low-cost sensors, automated machine learning, and big human mobility data. \textit{Frontiers in Environmental Science} \\(DOI:10.3389/fenvs.2023.1223160)\\
\years{2023}Liang, Y., \textbf{Yin, J.}, Park, S., Pan, B., Chi, G., and Miller, Z. (2023), Using social media user profiles to identify visitor demographics and origins in Yellowstone national park. \textit{Journal of Outdoor Recreation and Tourism} (DOI:10.1016/j.jort.2023.100620)\\
\years{2022}\textbf{Yin, J.}, Gao, Y. and Chi, G. (2022), An evaluation of geo-located Twitter data for measuring human migration. \textit{International Journal of Geographical Information Science}, 36(9), pp. 1830-1852, DOI: 10.1080/13658816.2022.2075878\\
\years{2022}Li, C., Zhao, J., \textbf{Yin, J.}  and Chi, G. (2022), Park access affects physical activity: new evidence from geolocated Twitter data analysis. \textit{Journal of Urban Design}\\(DOI: 10.1080/13574809.2022.2118698)\\
\years{2022}Alba, C., Pan, B., \textbf{Yin, J.},  Rice, W., Lin, M.,  Liang, Y. and Mitra, P. (2022), COVID-19's impact on visitation behavior to US national parks from communities of color - evidence from mobile phone data. \textit{Scientific Reports}, 12:13398, DOI: 10.1038/s41598-022-16330-z\\
\years{2022}Liang, Y., \textbf{Yin, J.},  Pan, B.,  Lin, M., Miller, L., Taff, D., Chi, G. (2022), Assessing the Validity of Mobile Device Data for Estimating Visitor Demographics and Visitation Patterns in Yellowstone National Park. \textit{Journal of Environmental Management}, 317(115410), DOI:10.1016/j.jenvman.2022.115410\\
\years{2022}Chi, G., \textbf{Yin, J.}, Morris, L., and Bodovski, Y., Global Tweet Mentions of COVID-19 (2022). \textit{Health Data Science}\\
\years{2021}\textbf{Yin, J.} and  Chi, G., Characterizing People’s Daily Activity Patterns in the Urban Environment: A mobility network approach with geographic context-aware Twitter data. \textit{Annals of the American Association of Geographers}, 111(7), pp.1967-1987\\
\years{2018}Gao, Y., Wang, S. Padmanabhan, A., \textbf{Yin, J.} and Cao, G. (2018). Mapping Spatiotemporal Patterns of Events Using Social Media: A Case Study of Influenza Trends. \textit{International Journal of Geographical Information Science}, 32(3), pp. 425-449\\
\years{2017}\textbf{Yin, J.}, Soliman, A., Yin, D. and Wang, S. (2017). Delineate Urban Boundaries in Great Britain from the Network of Large Scale Twitter User Spatial Interactions. \textit{International Journal of Geographical Information Science}, 31(7), pp. 1293-1313\\
\years{2017}Soliman, A., Soltani, Q., \textbf{Yin, J.}, Padmanabhan, A., and Wang, S (2017). Social sensing of urban land use based on analysis of Twitter users' mobility patterns. \textit{PLoS ONE}, 12(7): e0181657. DOI:10.1371/journal.pone.0181657\\
\years{2009}Jiang, B., \textbf{Yin, J.} and Zhao, S. (2009). Characterizing the human mobility pattern in a large street network. \textit{Physical Review E}, 80(2), 021136

\subsection*{Peer Reviewed Conferences}
\years{2023}\textbf{Yin, J.} and Chi, G. (2023). Fine-grained Population Estimates: Infusing Deep Learning into Dasymetric Mapping of Buildings. \textit{2023 North American Meetings of the Regional Science Association International}, San Diego, CA, November 15 -18, 2023\\
\years{2018}\textbf{Yin, J.}, Chi, G. and Van Hook, J. (2018). Evaluating the Representativeness in the Geographic Distribution of Twitter User Population. \textit{The 12th Workshop on Geographic Information Retrieval in ACM SIGSPATIAL 2018}, November 6-9, 2018, Seattle, Washington, USA

\subsection*{Data}
\years{2020}\textbf{Yin, J.}. Estimation of Twitter user demographics in the USA, 2014. \textit{Harvard Dataverse}, doi:10.7910/DVN/PKKAPK

\section*{Research Grants (selected)}
\textbf{Co-PI}. ``RAPID: Using Mobile Phone Data to Understand the Impacts of the COVID-19 Pandemic on Food Assistance Use in Alaska''. NSF (National Science Foundation) SOC grant (Award \# 2207436), $\$$200,000. Project dates: 02/01/2022--01/31/2025\\
\textbf{Co-PI}. ``RR: The Generalizability and Replicability of Twitter Data for Population Research''. NSF (National Science Foundation) SOC grant (Award \# SES-1823633), $\$$500,000. Project dates: 07/15/2018--06/30/2022\\
\textbf{Senior Personnel}. ``Pursuing Opportunities for Long-term Arctic Resilience for Infrastructure and Society''. NSF (National Science Foundation) NNA grant (Award \# NNA-1927827), $\$$3,000,000. Project dates: 01/01/2020--08/31/2025\\
\textbf{PI}. ``Mobility Patterns and Their Influence on Food Security in Alaskan Communities: Using Digital Trace Data for Social and Behavioral Research during the COVID-19 Pandemic''. SSRI COVID-19 Pandemic Emergency Research Funding, $\$$30,000 Project dates: 04/01/2020--03/30/2022\\
\textbf{PI}. ``A Deep Learning-enabled Dasymetric Mapping Approach for Fine-Scale Population Estimation in Alaska''. ICDS computational and Data Sciences seed grant award, $\$$35,000 Project dates: 05/01/2020--03/30/2022\\
\textbf{Co-PI}. ``The Generalizability and Replicability of Twitter Data for Population Research''. ICS seed grant award, $\$$35,000. Project dates: 07/01/2018--06/30/2019\\
\textbf{Co-PI}. ``Are we more willing to speak up through mobile phones?  A comparison of desktop vs. mobile political sharing on Facebook''. SSRC (Social Science Research Council) grant, $\$$49,017. Project dates: 07/01/2020--06/30/2022\\
\textbf{PI}. ``Deep learning for Fine-grained Population Estimation in the Arctic Region''. Nvidia Academic Hardware Grant Program. Project dates: 06/01/2023--05/30/2024\\
\textbf{PI}. ``A cloud computing enabled GIS platform for the integration and synthesis of multi-layer geospatial data sources in urban studies: Understanding urban dynamics from geospatial Big Data''. Microsoft Azure Data Science Research Award, $\$20,000$. Project dates: 12/04/2016--12/03/2017\\
\textbf{PI}. ``Data Science: Mining sequential mobility patterns from semantic Twitter user trajectories''. Microsoft Azure Data Science Research Award, $\$10,000$. Project dates: 05/01/2018--04/30/2019\\
\textbf{PI}. ``Mining Twitter User Demographics as a First-Step in Big Data for Population Research''. XSEDE (Extreme Science and Engineering Discovery Environment) Startup computational resource allocation award, 50,000 SUs $\&$ 40 TB storage (estimated $\$$2,794.78). Project dates: 03/27/2017--03/26/2019

\section*{Synergistic Activities}
\begin{enumerate}
\setlength\itemsep{0em}
\item Holds the complete SafeGraph Places and Monthly Patterns archive for 2018--2023 and has used it to fuse mobile phone and point-of-interest data, including tracking food assistance use in Alaska during COVID-19 (NSF RAPID 2207436) and validating mobile device data against visitor counts in Yellowstone National Park.
\item Developed methods for assessing the representativeness of social media data for population research and released estimated demographics of U.S. Twitter users as a public dataset (Harvard Dataverse, doi:10.7910/DVN/PKKAPK).
\item Developed deep learning and dasymetric mapping methods for fine-scale population estimation from building footprints and imagery (ICDS seed grant; Nvidia Academic Hardware Grant).
\item Led education and outreach at the CyberGIS Commons, University of Illinois at Urbana-Champaign (2014--2016), including short courses on Hadoop, Spark, and parallel databases for geospatial big data.
\item Mentors data science students at The George Washington University from capstone projects to research publications, including a first-author manuscript on machine-learning prediction of high-resolution PM2.5 (submitted to \textit{Environmental Data Science}).
\end{enumerate}

\section*{Best Paper Awards}
\years{2022}ENTER 2022 eTourism Conference, for Pan, B., Savanapelli, V., Shukla, A. and \textbf{Yin, J.} (corresponding author), Monitoring Human-Wildlife Interactions in National Parks with Crowdsourced Data and Deep Learning\\
\years{2013}13th International Conference on Computational Science and Its Applications (ICCSA 2013), for Truong-Hong, L., Thi, T.T.P., \textbf{Yin, J.} and Carswell, J.D., Detailed 3D Building Models for Google Earth Integration\\
\years{2012}Web and Wireless Geographical Information Systems (W2GIS 2012), for \textbf{Yin, J.} and Carswell, J.D., Effects of Variations in 3D Spatial Search Techniques on Mobile Query Speed vs Accuracy

\section*{Teaching (The George Washington University)}
\textbf{Graduate}: Visualization of Complex Data (DATS 6401); Data Science Capstone (DATS 6501)\\
\textbf{Undergraduate}: Data Warehousing (DATS 2104); Data Visualization for Data Science (DATS 2102)

\section*{Professional Service (selected)}
\textbf{Editorial}: Topic Editor, \textit{Geosciences}; Editorial Board Member, \textit{Journal of Geography and Cartography}; Guest Editor, \textit{Spatial Demography} special issue on Big Spatial Data in Demography (2023)\\
\textbf{Review panels}: National Science Foundation, Methodology, Measurement, and Statistics Program (2019)\\
\textbf{Journal review}: more than 50 journals, including \textit{Nature Communications}, \textit{Annals of the American Association of Geographers}, \textit{International Journal of Geographical Information Science}, and \textit{Applied Geography}\\
\textbf{University}: Research Advisory Council and GPAC Committee, Columbian College of Arts \& Sciences, The George Washington University (2026--)

\end{document}
